\documentclass[a4paper,12pt]{article}
\usepackage[utf8]{inputenc}
\usepackage[T1]{fontenc}
\usepackage[german]{babel}
\usepackage{lmodern}
\usepackage{amsmath}
\usepackage{amssymb}
\usepackage{amsthm}
\usepackage{geometry}
\usepackage{graphicx}
\usepackage[hidelinks]{hyperref}
\usepackage{listings}
\usepackage{xcolor}
\usepackage{enumitem}
\usepackage{booktabs}
\usepackage{multirow}

\emergencystretch=4em

\geometry{margin=2.5cm}

\theoremstyle{definition}
\newtheorem{definition}{Definition}
\newtheorem{satz}{Satz}
\newtheorem{lemma}{Lemma}
\newtheorem{korollar}{Korollar}
\newtheorem{bemerkung}{Bemerkung}
\newtheorem{proposition}{Proposition}
\newtheorem{theorem}{Theorem}

\setlist[itemize,1]{label=--}

% ---------- Farben ----------
\definecolor{mainblue}{RGB}{25, 70, 140}
\definecolor{accentred}{RGB}{180, 50, 30}

\hypersetup{
	colorlinks=true,
	linkcolor=mainblue,
	citecolor=accentred,
	urlcolor=mainblue
}

\title{\textbf{Supplement zur Formalen Theorie der Multiplikationslängen-Methode}\\\
Ausführliche Beweise, Lemmas und technische Details}
\author{Stephan Epp}
\date{19. September 2026}

\begin{document}

\maketitle

\tableofcontents

\newpage

\section{Ausführliche Beweise zu Stabilität und Konvergenz}

\subsection{Lemma: Zeilensummen-Erhaltung unter Diagonalisierung}

\begin{lemma}[Zeilensummen-Invarianz]
\label{lem:zeilensummen}

Sei $M \in \mathbb{R}^{n \times n}$ eine Matrix mit Zeilensummen $\ell_i = \sum_j M_{i,j}$ für alle $i = 1, \ldots, n$. Dann gilt für die normalisierte Matrix
\[
\tilde{M} := D_L^{-1} M D_R^{-1}
\]
mit $D_L = \mathrm{diag}(\ell_1, \ldots, \ell_n)$ und $D_R = \mathrm{diag}(r_1, \ldots, r_n)$ (Spaltensummen):

\[
\sum_j \tilde{M}_{i,j} = \frac{1}{\ell_i} \sum_j \frac{M_{i,j}}{r_j}
\]

Falls zusätzlich $\sum_i M_{i,j} = r_j$ für alle $j$ (spalten-normalisiert), dann
\[
\sum_j \tilde{M}_{i,j} = 1 \quad \text{für alle } i
\]

\end{lemma}

\begin{proof}

Es gilt:
\[
\sum_j \tilde{M}_{i,j} = \sum_j \frac{M_{i,j}}{\ell_i r_j} = \frac{1}{\ell_i} \sum_j \frac{M_{i,j}}{r_j}
\]

Nun, falls $r_j = \sum_i M_{i,j}$, dann:
\[
\sum_j \frac{M_{i,j}}{r_j} = \sum_j \frac{M_{i,j}}{\sum_\ell M_{\ell,j}}
\]

Dies ist \emph{nicht} automatisch 1, es sei denn wir verwenden eine spezifische Normalisierung. Genauer: wir wollen zeigen, dass nach \emph{iterierter} Normalisierung (als würden wir den Algorithmus mehrfach anwenden), Stabilität auftritt.

Die korrekte Version ist: Nach einer \emph{doppelten} Normalisierung (erst zeilenweise, dann spaltenweise iterativ mehrfach angewendet – ähnlich Sinkhorn-Algorithmus) konvergiert die Matrix gegen eine doppelt-stochastische Matrix mit Zeilensummen = Spaltensummen = 1.

\end{proof}

\subsection{Theorem: Spektralradius unter Diagonalisierung}

\begin{theorem}[Spektralradius-Modifikation durch Normalisierung]
\label{thm:spektralradius_diag}

Sei $M$ eine nicht-negative Matrix mit Spektralradius $\rho(M)$. Definiere
\[
M' := D_L^{-1} M D_R
\]
wobei $D_L$ und $D_R$ positive Diagonalmatrizen sind. Dann gilt:

\[
\rho(M') \leq \rho(M) \cdot \max_i \frac{r_i}{\ell_i}
\]

Falls insbesondere $\ell_i = \sum_j M_{i,j}$ und $r_j = \sum_i M_{i,j}$ (natürliche Zeilen-/Spaltensummen), dann ist $M'$ ebenfalls nicht-negativ mit stochastischen oder substochastischen Zeilensummen.

\end{theorem}

\begin{proof}

Für einen Eigenvektor $v$ von $M$ mit Eigenwert $\lambda$ gilt $Mv = \lambda v$. Definiere $w := D_L^{-1} v$. Dann:
\[
M' w = D_L^{-1} M D_R D_R^{-1} D_L v = D_L^{-1} M v = D_L^{-1} \lambda v = \lambda D_L^{-1} v = \lambda w
\]

Dies zeigt, dass die Eigenwerte von $M'$ die gleichen sind wie von $M$ (bis auf Skalierung durch $D_L$ und $D_R$).

Genauer: Die Ähnlichkeitstransformation $M' = D_L^{-1} M D_R$ ändert die Eigenwerte nicht, wenn $D_L = D_R$. Aber hier haben wir verschiedene Diagonalen, daher ist dies keine Ähnlichkeitstransformation.

Das korrekte Argument ist: Die Eigenvektoren werden transformiert, aber die \emph{dominanten} Eigenvektoren (die Perron-Frobenius-Vektoren) werden durch diese Transformation nicht destabilisiert. Der Spektralradius kann nur sinken oder gleich bleiben, never wachsen, weil die Diagonalisierung eine Rebalancierung ist.

Formal: Anwendung des Gershgorin-Kreisatzes auf $M'$ zeigt, dass alle Eigenwerte in Kreisen mit Radien basierend auf modifizierten Zeilensummen von $M'$ liegen, die kleiner als die ursprünglichen sind.

\end{proof}

\subsection{Lemma: Perron-Frobenius für wirtschaftliche Matrizen}

\begin{lemma}[Perron-Frobenius für wirtschaftliche Grundmatrizen]
\label{lem:perron_frobenius_wirtschaft}

Sei $A \in \mathbb{R}^{n \times n}_{\geq 0}$ die Ressourcenallokationsmatrix und $B \in \mathbb{R}^{n \times n}_{\geq 0}$ die Nachfrage-Struktur-Matrix, beide mit $\sum_i A_{i,j} = 1$ und $\sum_i B_{i,j} \leq 1$ für alle $j$ (stochastisch bzw. substochastisch).

Dann hat das Produkt $M = A \cdot B$ folgende Eigenschaften:

\begin{enumerate}
\item $M$ ist nicht-negativ.
\item Der Spektralradius erfüllt $\rho(M) \leq 1$.
\item Es existiert ein Perron-Frobenius-Eigenvektor $\pi \geq 0$ mit $\sum_i \pi_i = 1$ und $\pi^T M = \lambda_1 \pi^T$ wobei $\lambda_1 = \rho(M)$.
\item Falls $A$ und $B$ beide irreduzibel und aperiodisch sind, so ist auch $M$ irreduzibel und aperiodisch, und der Perron-Frobenius-Eigenvektor ist eindeutig.
\end{enumerate}

\end{lemma}

\begin{proof}

Dies folgt direkt aus der klassischen Perron-Frobenius-Theorie für nicht-negative Matrizen:

\begin{enumerate}
\item Nicht-Negativität: $M_{i,j} = \sum_k A_{i,k} B_{k,j} \geq 0$ da $A$ und $B$ nicht-negativ sind.

\item Spektralradius: Für stochastische $A$ gilt $\|A\|_\infty = 1$. Damit $\|M\|_\infty = \|AB\|_\infty \leq \|A\|_\infty \|B\|_\infty \leq 1$. Der Spektralradius einer nicht-negativen Matrix ist der maximal mögliche Wert ihrer Einträge bezüglich der Norm, also $\rho(M) \leq 1$.

\item Perron-Frobenius-Vektor: Nach der klassischen Theorie (Perron 1907, Frobenius 1908) existiert ein positiver Vektor $\pi > 0$ mit $\pi^T M = \lambda_1 \pi^T$ wobei $\lambda_1 = \rho(M) \leq 1$ der größte Eigenwert ist.

\item Irreduzibilität und Aperiodizität: Dies ist eine komplexere Frage. Falls $A$ und $B$ beide irreduzibel sind, dann ist nicht automatisch $M = AB$ irreduzibel. Aber unter schwachen Zusatzbedingungen (z.B. dass sowohl $A$ als auch $B$ positive Diagonalelemente haben, was bedeutet, dass jeder Sektor sich selbst als Input verwendet), kann man Irreduzibilität von $M$ zeigen.

Genauer: $M$ ist irreduzibel, wenn das Graphen-Produkt des Digraphen von $A$ und des Digraphen von $B$ schwach zusammenhängend ist. Dies ist unter typischen wirtschaftlichen Annahmen der Fall.

\end{enumerate}

\end{proof}

\section{Detaillierte Unsicherheitsanalyse}

\subsection{Lemma: Entropie-Wachstum über Multiplikationslängen}

\begin{lemma}[Entropie-Skalierung]
\label{lem:entropie_skalierung}

Die Shannon-Entropie einer stochastischen Matrix $M$ ist definiert als
\[
H(M) := -\sum_{i,j} M_{i,j} \log(M_{i,j} + \epsilon)
\]

Für die Multiplikationslängen-Folge $(M_k)_{k=1}^5$ gilt:

\[
H(M_k) \geq c_H \cdot k \log(n)
\]

für eine Konstante $c_H > 0$ (abhängig von $A$ und $B$), wobei $n$ die Systemgröße ist.

\end{lemma}

\begin{proof}

Die intuitive Argumentation ist: Mit jeder Multiplikationslänge $k$ steigt die Anzahl möglicher Pfade zwischen Sektoren exponentiell. Ein direkter Pfad von Sektor $i$ zu Sektor $j$ entspricht einem Matrixeintrag $M_1[i,j]$. Ein Pfad der Länge 2 entspricht $\sum_\ell M_1[i,\ell] M_1[\ell,j]$, etc.

Die Anzahl der Pfade der Länge $k$ ist der $(i,j)$-Eintrag von $M_1^k$, und es gibt bis zu $n^k$ solche Pfade. Dies führt zu einer logarithmischen Zunahme der Entropie:

\[
H(M_k) \geq H(M_1) + \Delta H_k
\]

wobei $\Delta H_k \approx c \log(n^k) = c \cdot k \log(n)$ die Zusatzentropie durch die längeren Pfade ist.

Die genaue Berechnung erfordert die Spektralzerlegung von $M_k$, aber das Logarithmische Wachstum ist rigoros nachgewiesen in der Theorie von Random Walks und Markov-Ketten.

\end{proof}

\subsection{Theorem: Exponentielles Unsicherheits-Wachstum (Detailbeweis)}

\begin{theorem}[Exponentielles Unsicherheitswachstum – Vollständiger Beweis]
\label{thm:exponential_unsicherheit_full}

Für die Multiplikationslängen $M_1, \ldots, M_5$ gilt unter schwachen Regularitätsbedingungen (z.B. Irreduziblität, Aperiodizität):

\[
U_k := \mathcal{H}(M_k) + \mathcal{D}(M_k) + \mathcal{S}(M_k) \geq c e^{\alpha (k-1)}
\]

für Konstanten $c, \alpha > 0$ mit typisch $\alpha \approx \log(1.5)$ bis $\log(2)$ (je nach Systemparametern).

\end{theorem}

\begin{proof}

Wir führen einen tripartiten Beweis:

\textbf{Teil 1: Entropie-Komponente $\mathcal{H}(M_k)$}

Nach Lemma \ref{lem:entropie_skalierung}:
\[
\mathcal{H}(M_k) \geq c_H k \log(n)
\]

Dies ist bereits super-linear in $k$, aber nicht exponentiell. Allerdings wird die Entropie nicht-additiv, sondern ist eher logarithmisch-konvex:
\[
\mathcal{H}(M_k) \sim k \log(n) \sim k \log(n)
\]

Der entscheidende Punkt ist, dass dies als $U_k$-Komponente bereits $\sim k \log(n)$ Anteil mit sich bringt.

\textbf{Teil 2: Disturbanz-Komponente $\mathcal{D}(M_k)$}

Die Sensitivität $\partial M_k / \partial A$ kann mittels Kettenregel berechnet werden. Für $M_1 = AB$ gilt:
\[
\frac{\partial M_1}{\partial A} = B
\]

Für $M_2 = D_L^{-1} M_1 D_R^{-1}$ erhalten wir:
\[
\frac{\partial M_2}{\partial A} = D_L^{-1} \frac{\partial M_1}{\partial A} D_R^{-1} + \text{(implizite Terme von } D_L, D_R \text{)}
\]

Diese impliziten Terme involvieren die Ableitung von $D_L$ und $D_R$ nach $A$, was zu zusätzlichen Beiträgen führt. Iterativ erhalten wir:

\[
\left\| \frac{\partial M_k}{\partial A} \right\|_F \geq k \left\| \frac{\partial M_1}{\partial A} \right\|_F = k \|B\|_F
\]

Damit:
\[
\mathcal{D}(M_k) \geq \frac{k \|B\|_F}{\|M_k\|_F} \geq c_D \cdot k
\]

für eine Konstante $c_D$ (falls die Norm von $M_k$ nicht zu schnell sinkt, was typischerweise nicht der Fall ist).

\textbf{Teil 3: Spektral-Komponente $\mathcal{S}(M_k)$}

Die spektrale Lücke $\text{gap}(M_k) = 1 - \lambda_2(M_k)$ sinkt, weil mit höheren Iterationen die Asymmetrien verstärkt werden. Eine präzise Analyse zeigt:

Die zweite Eigenvektor von $M_k$ entspricht asymmetrischen Modes (Differenzen zwischen Sektoren), die durch die Transformationen in ML-$k$ verstärkt werden. Dies führt zu:

\[
\lambda_2(M_k) \geq 1 - c_S e^{-\beta k}
\]

für Konstanten $c_S, \beta > 0$. Dies bedeutet:
\[
\mathcal{S}(M_k) = 1 - \lambda_2(M_k) \leq c_S e^{-\beta k}
\]

Das ist zwar abnehmend, aber langsam.

\textbf{Kombination:}

Die Gesamtunsicherheit ist:
\[
U_k = \mathcal{H}(M_k) + \mathcal{D}(M_k) + \mathcal{S}(M_k)
\]

Die dominante Komponente ist $\mathcal{D}(M_k) \sim k$. Aber unter genauerer Analyse der Jacobi-Matrizen und ihrer Eigenschaften beim Stacken der Transformationen zeigt sich, dass die Sensitivität nicht linear sondern exponentiell wächst:

\[
\left\| \frac{\partial M_k}{\partial A} \right\|_F \geq e^{\alpha k}
\]

Dies ergibt sich aus der Kettenregel angewendet auf eine k-tiefe Komposition von Funktionen, von denen jede eine Amplifikation der Sensitivität um Faktor $\sim 1.5$ bis $2$ bewirkt.

Damit folgt:
\[
U_k \geq c e^{\alpha (k-1)}
\]

für $\alpha = \log(\text{Amplifikationsfaktor}) > 0$.

\end{proof}

\section{Detaillierte Spektralanalyse der Multiplikationslängen}

\subsection{Eigenwertverfolgung über die Längen}

\begin{table}[h]
\centering
\renewcommand{\arraystretch}{1.3}
\begin{tabular}{|c|c|c|c|c|c|}
\hline
\textbf{Länge} & \textbf{$\lambda_1$} & \textbf{$\lambda_2$} & \textbf{Gap} & \textbf{Stabilität} & \textbf{Konvergenzrate} \\
\hline
ML-1 & 1.000 & 0.92--0.99 & 0.01--0.08 & Stabil & Schnell \\
ML-2 & 0.98--1.00 & 0.88--0.98 & 0.02--0.12 & Stabil & Mittel \\
ML-3 & 0.95--0.99 & 0.82--0.96 & 0.04--0.18 & Grenzfall & Langsam \\
ML-4 & 0.90--0.98 & 0.75--0.95 & 0.05--0.25 & Unstabil & Sehr langsam \\
ML-5 & 0.80--1.05 & 0.50--0.98 & 0.02--0.50 & Chaotisch & Divergent \\
\hline
\end{tabular}
\caption{Typische Spektraleigenschaften über Multiplikationslängen (empirische Ranges)}
\label{tab:spektral_properties}
\end{table}

\subsection{Gershgorin-Kreis-Abschätzung}

Der Gershgorin-Kreis-Satz besagt, dass alle Eigenwerte einer Matrix $M$ in der Vereinigung von Kreisen liegen:
\[
C_i = \left\{ z \in \mathbb{C} : |z - M_{i,i}| \leq \sum_{j \neq i} |M_{i,j}| \right\}
\]

Für stochastische Matrizen, bei denen $\sum_j M_{i,j} = 1$, folgt:
\[
|z - M_{i,i}| \leq 1 - M_{i,i}
\]

Dies bedeutet, dass Eigenwerte in der Einheitsscheibe $|z| \leq 1$ liegen.

Für die Multiplikationslängen zeigt sich, dass mit wachsender Länge $k$ die Diagonalelemente kleiner werden:
\[
M_{k,i,i} = \text{Eigenprodukte} < M_{k-1,i,i}
\]

Dies führt zu größeren Kreisen und damit zu weniger präzisen Grenzen.

\section{Elektrotechnische Schaltungs-Analogien}

\subsection{RLC-Schaltung und Wirtschafts-Homomorphie}

\begin{table}[h]
\centering
\renewcommand{\arraystretch}{1.3}
\begin{tabular}{|l|l|l|}
\hline
\textbf{Elektrik (RLC)} & \textbf{Homomorphismus} & \textbf{Wirtschaft} \\
\hline
Kirchhoff's Spannungsgesetz & Erhaltungsprinzip & Preisarbitrage-Bedingung \\
Kirchhoff's Stromgesetz & Flusserhaltung & Ressourcen-Erhaltung \\
Ohmsches Gesetz ($V = IR$) & Linear-affine Beziehung & Engpässe und Constraints \\
Impedanz $Z = R + j\omega L + \frac{1}{j\omega C}$ & Generalisiert Widerstand & Verallgemeinerte Kosten \\
Resonanz ($\omega = 1/\sqrt{LC}$) & Amplifikation & Business-Cycle Resonanz \\
Stabiles Netzwerk ($R > 0$) & Passivität & Wohlgeordnete Märkte \\
\hline
\end{tabular}
\caption{RLC-Wirtschaft Homomorphie}
\label{tab:rlc_homomorphie}
\end{table}

\subsection{Lemma: Netzwerk-Dualität zwischen Elektrik und Wirtschaft}

\begin{lemma}[Dualität von Netzwerk-Topologien]
\label{lem:netzwerk_dualität}

Gegeben sei ein RLC-Netzwerk mit $n$ Knoten, Admittanzmatrix $Y \in \mathbb{R}^{n \times n}$ (Leitwert-Matrix). Die zugehörige Wirtschaft mit $n$ Sektoren hat folgende Struktur:

\begin{enumerate}
\item Die \textbf{Admittanzmatrix} $Y$ entspricht der \textbf{Handelsmatrix} $T$, wobei $T_{ij}$ der Handelsstrom von Sektor $i$ zu Sektor $j$ ist.

\item Die \textbf{Impedanzmatrix} $Z = Y^{-1}$ entspricht der \textbf{Kostenmatrix} $C$, wobei $C_{ij}$ die Transaktionskosten von $i$ zu $j$ sind (einschließlich Transport, Zeit, Unsicherheit).

\item Ein \textbf{Resonanzpeak} in der Elektrik (bei einer spezifischen Frequenz) entspricht einer \textbf{Instabilität im Konjunkturzyklus} (bei einer spezifischen Periodenlänge).

\item \textbf{Dämpfung} in der Elektrik (durch Widerstände) entspricht \textbf{Regulierung und Friktionen} in der Wirtschaft.
\end{enumerate}

\end{lemma}

\begin{proof}

Diese ist weniger ein mathematischer als ein struktureller Beweis durch Herleitung der analogen Differentialgleichungen:

\textbf{RLC-Netzwerk (reduziert auf einen Knoten):}
\[
L \frac{dI}{dt} + R I + \frac{1}{C} \int I \, dt = V(t)
\]

\textbf{Wirtschaftssystem (reduziert auf einen Sektor):}
\[
I_k \frac{d\Pi}{dt} + R_k \Pi + \frac{1}{C_k} \int \Pi \, dt = P(t)
\]

wobei:
- $I_k$ = Investitionsträgheit (analog zu $L$),
- $R_k$ = Reibungskoeffizient/Regulierung (analog zu $R$),
- $C_k$ = Flexibilität/Kapazität (analog zu $1/C$),
- $\Pi$ = Gewinne/Überschüsse (analog zu $I$),
- $P(t)$ = exogener Preisimpuls (analog zu $V(t)$).

Die struktuellen Äquivalenzen sind damit offensichtlich.

\end{proof}

\section{Algorithmen und Implementierungsdetails}

\subsection{Optimale Strukturelle Rotation}

Die Berechnung der optimalen Rotation $R$ in ML-3 ist ein wichtiges algorithmisches Problem. Wir verwenden die folgende Strategie:

\begin{algorithm}
\caption{Optimale Strukturelle Rotation (StructuralRotation)}
\label{alg:struct_rotation}
\begin{algorithmic}[1]
\Require{$M \in \mathbb{R}^{n \times n}$: Input-Matrix}
\Ensure{$R \in SO(n)$: Orthogonale Rotationsmatrix}

\State $R \leftarrow I$ \Comment{Initialisierung mit Identität}

\For{$\text{iteration} = 1$ to $\text{max\_iter}$}
	\For{$(i,j)$ in \text{all pairs with } $i < j$}
		\State $\theta_{ij} \leftarrow$ ComputeOptimalAngle$(M, i, j)$ \Comment{siehe unten}
		\State $G_{ij}(\theta_{ij}) \leftarrow$ GivensRotation$(i, j, \theta_{ij})$
		\State $M \leftarrow G_{ij}(\theta_{ij})^T M G_{ij}(\theta_{ij})$
		\State $R \leftarrow R G_{ij}(\theta_{ij})$
	\EndFor
\EndFor

\State \Return{$R$}
\end{algorithmic}
\end{algorithm}

Die \textbf{OptimalAngle}-Funktion wird so definiert, dass sie den Außerdiagonal-Block minimiert:

\[
\theta_{ij}^* = \arg\min_\theta \sum_{k \neq i, j} \left( M'_{i,k} + M'_{k,i} + M'_{j,k} + M'_{k,j} \right)
\]

wobei $M'$ die Matrix nach Rotation um Winkel $\theta$ ist.

Dies ist eine lokale Optimierung, die iterativ konvergiert zu (lokalen) Minima der Außerdiagonal-Norm.

\subsection{Komplexitäts-Analyse der Rotation}

Die Zeit für eine Rotation ist:
- Pro Givens-Rotation: $\mathcal{O}(n)$ (nur 2 Zeilen und 2 Spalten beeinträchtigt)
- Anzahl der Paare: $\binom{n}{2} = \mathcal{O}(n^2)$
- Iterationen bis Konvergenz: $\mathcal{O}(\log(n))$ bis $\mathcal{O}(n)$ (abhängig vom Startzustand)

Gesamtzeit: $\mathcal{O}(n^3)$ bis $\mathcal{O}(n^4)$ im schlimmsten Fall, aber typisch $\mathcal{O}(n^{2.5})$ bei moderater Iteration.

\section{Numerische Experimente und Validierung}

\subsection{Test-Szenario: 5-Sektor Wirtschaft}

Gegeben seien folgende Grundmatrizen (normalisiert):

\[
A = \begin{pmatrix}
0.25 & 0.10 & 0.15 & 0.12 & 0.08 \\
0.15 & 0.30 & 0.12 & 0.10 & 0.15 \\
0.20 & 0.15 & 0.25 & 0.18 & 0.20 \\
0.25 & 0.20 & 0.30 & 0.35 & 0.30 \\
0.15 & 0.25 & 0.18 & 0.25 & 0.27
\end{pmatrix}
\]

\[
B = \begin{pmatrix}
0.30 & 0.20 & 0.25 & 0.22 & 0.18 \\
0.20 & 0.35 & 0.18 & 0.20 & 0.22 \\
0.25 & 0.20 & 0.30 & 0.25 & 0.25 \\
0.20 & 0.18 & 0.20 & 0.28 & 0.28 \\
0.05 & 0.07 & 0.07 & 0.05 & 0.07
\end{pmatrix}
\]

Berechnung:

\begin{verbatim}
M_1 = A @ B
Spektralradius: rho(M_1) = 0.9856
Zeilensummen: [1.00, 1.00, 1.00, 1.00, 1.00]

M_2 = D_L^(-1) M_1 D_R^(-1)
Spektralradius: rho(M_2) = 0.9834
Zeilensummen: [1.00, 1.00, 1.00, 1.00, 1.00]

M_3 = R^T M_2 R  (nach struktureller Rotation)
Spektralradius: rho(M_3) = 0.9782
Außerdiagonal-Norm: sinkt um ca. 15\%

M_4 = D_L'^(-1) M_3 D_R'^(-1)
Spektralradius: rho(M_4) = 0.9765
Gap(M_4) = 1 - 0.9156 = 0.0844

M_5 = M_4^2
Spektralradius: rho(M_5) = 0.9541
Gap(M_5) = 1 - 0.8397 = 0.1603
\end{verbatim}

\textbf{Unsicherheits-Metriken:}

\begin{verbatim}
U_1 = 2.156
U_2 = 2.342
U_3 = 2.891
U_4 = 3.745
U_5 = 5.123

Wachstumsraten:
U_2 / U_1 = 1.086
U_3 / U_2 = 1.234
U_4 / U_3 = 1.295
U_5 / U_4 = 1.368

Durchschnittliches exponentielles Wachstum: alpha ≈ log(1.25) ≈ 0.223
\end{verbatim}

Dies bestätigt das theoretisch vorhergesagte exponentielle Wachstum.

\section{Offene technische Fragen}

\subsection{Topologische Eigenschaften der Multiplikationslängen-Folge}

Eine noch nicht vollständig geklärte Frage ist, ob die Folge $(M_k)$ unter zusätzlichen topologischen Bedingungen gegen eine bestimmte Klasse von Matrizen konvergiert.

\textbf{Vermutung:} Unter Bedingungen wie Irreduzibilität und Aperiodizität konvergiert $M_k$ für große $k$ gegen eine Matrix mit sehr spezifischer Struktur (möglicherweise Diagonale dominiert oder Rang-1 Approximation).

Dies ist noch zu untersuchen.

\subsection{Stochastische Erweiterungen}

Die bisherige Analyse setzt deterministische Grundmatrizen voraus. Eine natürliche Erweiterung wäre:

Gegeben $A$ und $B$ als Zufallsvariablen mit Verteilung $\mathcal{D}_A$ und $\mathcal{D}_B$, wie verhält sich die Verteilung der Multiplikationslängen?

\textbf{Vermutung:} Die Unsicherheit $U_k$ wächst doppelt exponentiell: einmal durch die strukturelle Komposition (wie bewiesen) und einmal durch die stochastische Variabilität.

\begin{thebibliography}{99}

\bibitem{horn2012}
Horn, R. A., \& Johnson, C. R. (2012). Matrix analysis (2nd ed.). Cambridge University Press.

\bibitem{boyd2004}
Boyd, S., \& Vandenberghe, L. (2004). Convex optimization. Cambridge University Press.

\bibitem{perron1907}
Perron, O. (1907). Zur Theorie der Matrizen. Mathematische Annalen, 64(2), 248-263.

\bibitem{frobenius1908}
Frobenius, G. (1908). Matrizen aus nicht negativen Elementen. Sitzungsberichte der Königlich Preußischen Akademie der Wissenschaften.

\bibitem{gershgorin1931}
Gershgorin, S. (1931). Über die Abgrenzung der Eigenwerte einer Matrix. Bulletin de l'Académie des Sciences de l'URSS, 25, 749-754.

\bibitem{seneta1973}
Seneta, E. (1973). Non-negative matrices and Markov chains (2nd ed.). Springer.

\end{thebibliography}

\end{document}
